\chapter{Abstraction}

\section{Vai trò của Abstraction}
\defaultauthor

Abstraction che giấu chi tiết không cần thiết và cung cấp contract rõ ràng cho tầng phía trên. IMOVE dùng abstraction để tách UI, HTTP API, agent, service, typed model, context và persistence.

\section{Các tầng abstraction thực tế}
\defaultauthor

\begin{table}[H]
\centering
\begin{tabularx}{\textwidth}{p{3.5cm}X X}
\toprule
\textbf{Layer} & \textbf{Implementation} & \textbf{Chi tiết được che giấu} \\
\midrule
Frontend & Pages, components, hooks, API service & HTTP, auth header, local cache \\
Router & trips, places, alerts, transit, preferences, chat & Request validation và HTTP response \\
Agent & Planning, Adaptation, Memory, Chat & Use-case orchestration \\
Service & OneMap, LTA, OpenWeather, Gemini, Scoring & API protocol và thuật toán chuyên biệt \\
Model & Pydantic request/response models & Dictionary shape không kiểm soát \\
Context & \code{ContextSnapshot} & Cách lấy weather và tính peak hours \\
Persistence & Supabase, in-memory, localStorage & Chi tiết lưu trữ theo phạm vi \\
\bottomrule
\end{tabularx}
\caption{Các tầng abstraction trong IMOVEV2.}
\end{table}

\section{Router, Agent và Service Abstraction}
\defaultauthor

FastAPI routers là HTTP boundary. Agents thực hiện use case cấp cao, còn services đóng gói external API hoặc thuật toán chuyên biệt. Ví dụ, Planning Agent yêu cầu alternatives đã chuẩn hóa mà không cần expose chi tiết token OneMap; frontend nhận \code{TripPlan} mà không cần biết thuật toán greedy hoặc scoring.

\begin{figure}[H]
\centering
\begin{tikzpicture}[node distance=1.1cm]
  \node[imovenode] (front) {React Frontend};
  \node[imovenode,below=of front] (router) {FastAPI Routers};
  \node[imoveagent,below=of router] (agent) {Agents};
  \node[imoveagent,below=of agent] (service) {Services};
  \node[imoveexternal,below=of service] (external) {External APIs};
  \draw[imovearrow] (front) -- (router);
  \draw[imovearrow] (router) -- (agent);
  \draw[imovearrow] (agent) -- (service);
  \draw[imovearrow] (service) -- (external);
\end{tikzpicture}
\caption{Luồng abstraction từ giao diện đến external APIs.}
\label{fig:abstraction-layers}
\end{figure}

\section{Model và Context Abstraction}
\defaultauthor

\code{TripPlan}, \code{DayPlan}, \code{LegResponse} và \code{AlternativeRoute} tạo contract cho itinerary và route. \code{UserPreferenceProfile} và \code{ModeConstraints} đóng gói preference. \code{ContextSnapshot} gom lượng mưa và thời gian hiện tại, đồng thời cung cấp \code{rain_level} và \code{is_peak_hours} cho scoring.

\section{Data and Persistence Abstraction}
\defaultauthor

\begin{itemize}
  \item Curated \code{singapore_places.json} được load và validate khi backend khởi động.
  \item Supabase PostgreSQL lưu trips, route legs, places, alerts, feedback và preferences \cite{supabase}.
  \item Supabase Auth cung cấp identity cho dữ liệu người dùng.
  \item In-memory store hỗ trợ active session và fallback.
  \item Frontend localStorage lưu metadata/cache cục bộ.
\end{itemize}

\section{Hybrid AI-assisted Planning}
\defaultauthor

Code xác định chịu trách nhiệm validation, scheduling, routing, scoring và fallback. Gemini hỗ trợ tác vụ ngôn ngữ như resolve tên địa điểm, suggestion, warning và chatbot \cite{gemini}. Ranh giới này hạn chế hallucination ảnh hưởng đến route, chi phí hoặc thời gian.

\section{Giới hạn của Abstraction hiện tại}
\defaultauthor

Hệ thống chưa có interface chung cho nhiều AI provider, chưa phải distributed microservice architecture, một phần persistence vẫn nằm trong routers và curated dataset được Planning Agent load trực tiếp. Đây là các hướng refactor hợp lý nhưng không được mô tả như khả năng đã hoàn thành.
